\documentclass[11pt,a4paper]{article}

\usepackage[utf8]{inputenc}
\usepackage[T1]{fontenc}
\usepackage[english]{babel}
\usepackage{lmodern}
\usepackage{microtype}
\usepackage{geometry}
\usepackage{booktabs}
\usepackage{array}
\usepackage{amsmath,amssymb}
\usepackage{graphicx}
\usepackage{xcolor}
\usepackage{hyperref}
\usepackage{enumitem}
\usepackage{listings}
\usepackage{caption}

\geometry{margin=2.4cm}
\hypersetup{colorlinks=true,linkcolor=blue,citecolor=blue,urlcolor=blue}
\setlist{nosep}

\title{\textbf{ReversaFeynman: Evidence-Governed Software Engineering Agents for Reverse Documentation, Repair, and Adaptive Evaluation}}
\author{Marcelo Claro Laranjeira}
\date{September 2026}

\begin{document}
\maketitle

\begin{abstract}
Large-language-model-based software engineering agents increasingly operate over legacy repositories, execute tools, modify code, and learn from prior trajectories. These capabilities create a methodological tension: confidence, memory, benchmark performance, and learned policy scores can improve decision making, but none of them is equivalent to direct evidence about program behavior. This paper presents \emph{ReversaFeynman}, an evidence-governed extension of the Reversa reverse documentation engineering framework. The system preserves Reversa's multi-agent extraction of operational specifications while adding explicit epistemic states, Feynman-inspired falsifiability gates, governed human teach-back, adaptive policy evaluation, longitudinal Hermes memory and skill proposals, an evidence governor, and a Software Engineering Intelligence layer for code localization, sandboxable execution, multi-candidate repair, static and mutation testing, tracing, benchmarking, durable workflows, MCP-ready interoperability, and offline prompt/skill optimization. The implementation follows specification-driven development (SDD) and test-driven development (TDD), and keeps external frameworks optional through adapters. We define a reproducible evaluation protocol comparing the original Reversa, minimal repair baselines, and ReversaFeynman across specification quality, repair success, calibration, regressions, latency, cost, and epistemic error. We deliberately do not claim empirical superiority in this version; instead, we contribute an implemented architecture and an experimental protocol intended to make such claims testable.
\end{abstract}

\noindent\textbf{Keywords:} software engineering agents; reverse engineering; legacy systems; LLM agents; evidence governance; program repair; software benchmarks; calibration; SDD; TDD.

\section{Introduction}

Software engineering agents based on large language models (LLMs) are moving from isolated code generation toward repository-scale activities such as maintenance, repair, migration, testing, and documentation. This transition increases the importance of reliable context and explicit correctness criteria. Legacy repositories are especially difficult because business rules, architectural decisions, operational exceptions, and hidden contracts are often distributed across code, configuration, data, tests, and maintenance practices.

Reversa was proposed as a reverse documentation engineering framework that converts legacy software into traceable operational specifications for AI agents through a multi-agent pipeline and explicit confidence marking \cite{macedo2026reversa}. Its central contribution is not autonomous coding itself, but the construction of a specification layer that makes implicit system knowledge more explicit and traceable.

The present work starts from a complementary problem. Once specifications are available, an agentic system still needs to answer questions such as: Which claims were directly observed? Which were inferred? What evidence could refute a claim? Can memory or a learned policy be trusted as evidence? How should the system select between several candidate patches? How can a repair be evaluated without conflating benchmark performance with truth?

We introduce \emph{ReversaFeynman}, an evidence-governed extension of Reversa. The framework retains the original reverse documentation pipeline and adds five classes of capability: (i) epistemic governance through explicit evidence states and Feynman-inspired gates; (ii) governed memory, skill evolution, and evidence collection through a Hermes integration; (iii) adaptive routing and offline policy evaluation; (iv) operational software engineering intelligence for code localization, isolated execution, multi-candidate repair, static analysis, mutation testing, observability, and benchmarking; and (v) reproducible development of the framework itself through SDD and TDD.

The core design principle is simple but strict: \emph{learned confidence is not direct evidence}. In ReversaFeynman, policy confidence, memory, reward, benchmark score, mutation score, static-analysis rank, and human confidence may guide action selection, but they do not independently promote a claim to the strongest epistemic state. Promotion requires direct, traceable evidence governed by deterministic local rules.

This paper makes the following contributions:

\begin{enumerate}
  \item an evidence-governed architecture that extends reverse documentation engineering with explicit epistemic states and falsifiability-oriented gates;
  \item a software engineering intelligence layer that separates localization, candidate generation, execution, independent validation, and ranking;
  \item an adapter-oriented design that keeps execution runtimes, scanners, observability systems, orchestration frameworks, and optimizers optional;
  \item a reproducible evaluation protocol for comparing the original Reversa, minimal repair baselines, and ReversaFeynman on both engineering and epistemic metrics;
  \item an implementation discipline in which the framework's own evolution is specified and tested through SDD/TDD artifacts committed with the code.
\end{enumerate}

We do not claim that ReversaFeynman is empirically superior to Reversa or to state-of-the-art coding agents in this manuscript. Instead, we distinguish architectural capability from empirical superiority and define the experiments required to evaluate the latter.
\section{Background and Problem Formulation}

\subsection{Reverse documentation engineering}

The original Reversa framework addresses a practical mismatch between agentic software maintenance and legacy codebases: coding agents need explicit behavioral contracts, while mature systems often encode critical rules only implicitly. Reversa structures discovery as a sequence of specialized roles that map project structure, excavate implementation details, infer domain and architecture, synthesize specifications, and review unresolved gaps \cite{macedo2026reversa}. Confidence is represented explicitly so that uncertain extractions are not silently treated as confirmed facts.

ReversaFeynman preserves this foundation. The extension does not redefine reverse documentation engineering; it adds a governance layer for what happens after claims, specifications, candidate repairs, trajectories, and learned policies begin to interact.

\subsection{Epistemic failure modes in agentic software engineering}

Repository-scale agents can fail even when individual model outputs appear plausible. We focus on six recurrent failure modes:

\begin{enumerate}
  \item \textbf{label substitution}: a named pattern or architecture is repeated without a causal explanation;
  \item \textbf{provenance loss}: a strong claim is detached from the code, test, execution, log, dataset, or artifact that supports it;
  \item \textbf{observation--inference collapse}: an interpretation is represented as if it had been directly observed;
  \item \textbf{non-falsifiable validation}: a test suite passes, but no meaningful oracle exists that could reveal a relevant behavioral error;
  \item \textbf{cargo-cult repair}: an established pattern is copied despite insufficient evidence that it addresses the local requirement;
  \item \textbf{confidence substitution}: memory, model confidence, human fluency, reward, or benchmark score is treated as equivalent to truth.
\end{enumerate}

The design of ReversaFeynman follows a Popperian intuition: strong engineering claims should expose conditions under which they could be shown false \cite{popper1959logic}. This does not turn software development into scientific experimentation in a strict philosophical sense; rather, it motivates operational gates that demand provenance, counterexamples, and executable checks when available.

\subsection{Epistemic states}

The framework uses four system-level states:

\begin{center}
\begin{tabular}{p{3.0cm}p{10.5cm}}
\toprule
\textbf{State} & \textbf{Meaning} \\
\midrule
\texttt{OBSERVED} & Supported by direct, traceable evidence accepted by the evidence governor. \\
\texttt{INFERRED} & Reasoned from available evidence but not directly established. \\
\texttt{UNVERIFIED} & Plausible or proposed, but not yet verified. \\
\texttt{BLOCKED} & Verification cannot proceed because required evidence, access, or a decision is missing. \\
\bottomrule
\end{tabular}
\end{center}

Human source quality is modeled separately from system evidence. A domain specialist may be classified as \texttt{HUMAN-VALIDATED}, \texttt{HUMAN-PARTIAL}, or \texttt{HUMAN-CONFLICT}; none of these labels is automatically equivalent to \texttt{OBSERVED}. This separation is essential when human explanation is the best available source but still requires traceability or executable corroboration.
\section{Method}

\subsection{Specification-driven and test-driven evolution}

The framework is developed using a combined specification-driven development (SDD) and test-driven development (TDD) process. Each architectural evolution is first expressed as a versioned specification containing invariants and acceptance criteria. Tests are then added before or together with the implementation so that the contract is executable. The implementation is subsequently hardened against boundary cases, compatibility regressions, and accidental privilege escalation.

For the Software Engineering Intelligence v5 evolution, the specification defines twelve invariants, including optional external integrations, prohibition of evidence promotion from learned scores, localization-before-editing, support for multiple repair candidates, execution-provider decoupling, falsifiability-oriented quality gates, observational benchmarking, shadow-first optimization, durable checkpoints, and strict schemas without numeric coercion.

The development cycle is:

\begin{center}
\texttt{SPEC $\rightarrow$ RED $\rightarrow$ GREEN $\rightarrow$ REFACTOR/HARDEN $\rightarrow$ PR/CI}
\end{center}

This process is itself part of the research artifact. Architectural claims about safety boundaries are therefore represented both in prose specifications and in executable regression tests.

\subsection{Feynman-inspired evidence gates}

The evidence and understanding layer defines seven gates:

\begin{enumerate}
  \item \textbf{FEG-01: Name is not understanding.} The mechanism should be explainable beyond merely naming a pattern or component.
  \item \textbf{FEG-02: Evidence and provenance.} Strong claims require a traceable supporting source.
  \item \textbf{FEG-03: Observation is not inference.} Direct observations and interpretations must remain distinguishable.
  \item \textbf{FEG-04: Falsifiability.} When applicable, the system should identify an executable test, oracle, counterexample, or observation that could refute the claim.
  \item \textbf{FEG-05: Anti-cargo-cult.} A proposed technique must address a demonstrated local need rather than merely imitate common practice.
  \item \textbf{FEG-06: Minimum uncertainty-reducing experiment.} The system identifies the smallest additional check that materially reduces uncertainty.
  \item \textbf{FEG-07: Teach-back boundary.} When repository evidence is insufficient and a human expert is required, the protocol tests causal explanation and transfer to a variant scenario rather than verbal fluency.
\end{enumerate}

The first six gates form the base engineering audit. The seventh is conditional and applies when domain knowledge materially affects a requirement, risk, rule, or legacy interpretation.

\subsection{Evidence governance}

The Hermes Evidence Governor replaces an earlier standalone Evidence Guard as the active operational boundary for evidence transitions. Hermes may collect candidate evidence, but a remote runtime, memory item, user model, skill proposal, learned policy, or trust score cannot independently establish \texttt{OBSERVED}. The local governor accepts promotion only when direct evidence has a recognized kind and a non-empty traceable reference.

Accepted baseline evidence kinds include code, contracts, tests, execution results, logs, datasets, and artifacts. This makes the evidence boundary deterministic even when the surrounding orchestration is adaptive.

\subsection{Adaptive evaluation without epistemic privilege}

The adaptive layer records learning events, computes operational reward, detects drift, and evaluates shadow policies offline. Policy confidence is interpreted as an estimated probability of success rather than as evidence about a program claim. Calibration can therefore be assessed using proper scoring rules such as the Brier score \cite{brier1950verification}, while the evidence governor remains independent of those calibration results.

The runtime separates \texttt{decide()} from \texttt{observe()} to reduce look-ahead leakage: a recommendation is made using only prior history, the baseline workflow executes, and only then is the outcome recorded for future learning.
\section{Architecture}

\subsection{Layered organization}

ReversaFeynman is organized as a set of cooperating layers rather than a single monolithic agent. The original Reversa pipeline remains responsible for extracting operational specifications from legacy systems. The Feynman layer audits understanding and evidence. MCI-style routing and adaptive evaluation support action selection. Hermes provides longitudinal memory, skill proposals, trajectories, and evidence collection, while the Hermes Evidence Governor remains the operational authority for epistemic state transitions. The Software Engineering Intelligence v5 layer adds repository-scale coding capabilities.

The architecture can be summarized as:

\begin{center}
\begin{minipage}{0.88\linewidth}
\begin{verbatim}
Legacy / feature / issue
        |
Reversa Discovery + SDD
        |
Feynman / Evidence Governance
        |
Code Intelligence Graph
        |
Localization
        |
Repair Laboratory --> N patch candidates
        |
Execution Fabric
        |
Quality Gates: tests / static / mutation
        |
Hermes Evidence Governor
        |
Trace + ReversaBench
        |
MCI / OPE / Adaptive Governance
        |
Offline Optimizer (shadow only)
\end{verbatim}
\end{minipage}
\end{center}

\subsection{Code Intelligence Graph}

Repository-scale repair requires more than textual search. Inspired by repository maps based on Tree-sitter and graph ranking, the Code Intelligence layer normalizes symbols, files, and directed relations into a graph that can be ranked deterministically for context selection. External providers may use Tree-sitter, ast-grep, Semgrep, or other structural analyzers, but the core consumes a provider-neutral contract. A graph rank is a relevance signal only; it has no epistemic authority.

\subsection{Execution Fabric}

Agent logic is separated from execution infrastructure. The Execution Fabric accepts injected providers and normalizes exit code, standard output, standard error, artifacts, status, task identity, and provider identity. This boundary permits local execution, containers, or remote execution systems such as SWE-ReX without coupling the ReversaFeynman core to a particular runtime \cite{swerex2026}.

\subsection{Repair Laboratory}

The repair process follows the sequence repository $\rightarrow$ file $\rightarrow$ symbol $\rightarrow$ edit region $\rightarrow$ candidate patches $\rightarrow$ independent validation $\rightarrow$ ranking. This is influenced by the localization--repair--validation decomposition explored by Agentless \cite{xia2024agentless}. ReversaFeynman does not assume that the first patch generated by an LLM is correct. Multiple candidates may be scored using test success, static analysis, mutation score, and regression penalties.

\subsection{Quality and falsifiability gates}

Quality evaluation combines conventional tests, static analysis, and optional mutation testing. Static analyzers such as Semgrep can identify semantic or policy violations \cite{semgrep2026}. Mutation systems such as StrykerJS can alter program semantics and measure whether the current test oracle detects the change \cite{stryker2026}. In ReversaFeynman, a mutation score strengthens FEG-04 by probing whether a test suite can refute relevant behavioral changes; the score is diagnostic and never constitutes direct evidence by itself.

\subsection{Observability, benchmarking, and optimization}

Traces record spans, attributes, timing, and outcomes without changing decisions. They may be exported to observability systems such as Phoenix/OpenTelemetry \cite{phoenix2026}. ReversaBench stores comparable records for named variants and reports success rate, latency, cost, and regressions while explicitly marking the comparison as non-causal unless the experimental design supports causal identification.

Durable execution is exposed through injected checkpoint stores and can be mapped to stateful workflow runtimes such as LangGraph \cite{langgraph2026}. MCP-ready handlers expose allowlisted tools through a provider-neutral gateway \cite{mcp2026}. Finally, prompt or skill optimization is restricted to shadow proposals that require offline evaluation and review. This boundary is compatible with programmatic optimizers such as DSPy and reflective search methods such as GEPA \cite{khattab2024dspy,agrawal2026gepa}.
\section{Implementation}

The reference implementation is maintained in the same repository as the paper. The core is a Node.js package with versioned agent skills, installer support for multiple coding-agent harnesses, adaptive integration modules, Hermes integration modules, and the Software Engineering Intelligence v5 layer.

\subsection{Software Engineering Intelligence v5}

Version 5 is implemented as ten incremental capabilities:

\begin{enumerate}
  \item \textbf{Strict Contract Layer}: versioned contracts use JSON Schema 2020-12 semantics and reject silent numeric coercion. An optional Ajv adapter can compile strict schemas when Ajv is installed.
  \item \textbf{Code Intelligence Graph}: normalizes symbols, relations, and a deterministic relevance ranking; structural extraction is delegated to optional providers.
  \item \textbf{Execution Fabric}: normalizes execution outcomes while keeping Docker, SWE-ReX, OpenHands, local shells, and remote runtimes outside the mandatory core.
  \item \textbf{Repair Laboratory}: requires a traceable localization before candidate creation and supports ranking of multiple repairs.
  \item \textbf{Quality/Falsifiability Gates}: combines tests, static-analysis findings, mutation-testing results, and regression penalties.
  \item \textbf{Observability}: records non-authoritative execution traces and spans.
  \item \textbf{ReversaBench}: stores comparable variant records and reports engineering metrics with an explicit \texttt{causal\_claim=false} default.
  \item \textbf{MCP-ready Gateway}: dispatches only registered/allowlisted tools and is compatible with an external MCP transport layer.
  \item \textbf{Durable Workflow Adapter}: persists step completion through injected load/save functions and avoids re-executing completed steps after resume.
  \item \textbf{Offline Prompt/Skill Optimizer}: produces shadow proposals with \texttt{auto\_apply=false}, offline-evaluation requirements, and no evidence authority.
\end{enumerate}

\subsection{Optional dependency boundary}

A deliberate implementation choice is that external frameworks remain optional. OpenHands, SWE-ReX, Semgrep, StrykerJS, Phoenix, LangGraph, MCP SDKs, DSPy, and related systems are integrated through adapters, transports, or executable providers rather than mandatory runtime dependencies. This keeps the ReversaFeynman package small and makes it possible to compare providers experimentally.

\subsection{Internal safety contracts}

The following invariants are enforced by tests:

\begin{itemize}
  \item numeric strings are rejected where a numeric contract is required;
  \item unknown execution providers and non-allowlisted tool calls fail closed;
  \item repair candidates require repository localization metadata;
  \item benchmark, mutation, rank, trace, and optimization results carry \texttt{evidence\_authority=false};
  \item the optimizer is shadow-first and cannot auto-apply changes;
  \item completed durable-workflow steps are not executed twice after checkpoint restore;
  \item external frameworks do not become mandatory package dependencies;
  \item evidence transitions remain delegated to the Hermes Evidence Governor.
\end{itemize}

\subsection{Provenance and compatibility}

The implementation explicitly distinguishes inherited capabilities from derived extensions. Discovery, reverse documentation engineering, operational specifications, and the foundational multi-agent pipeline are attributed to the original Reversa work \cite{macedo2026reversa}. Feynman gates, adaptive governance, offline policy evaluation, Hermes bridges/governance, and Software Engineering Intelligence v5 are extensions in the ReversaFeynman line. The repository preserves compatibility shims where an internal component has been replaced, reducing migration risk for existing consumers.
\section{Evaluation Protocol}

This version of the paper reports an implemented framework and a pre-specified evaluation protocol rather than completed comparative benchmark results. The purpose of this section is to make future superiority claims falsifiable and reproducible.

\subsection{Research questions}

\begin{description}
  \item[RQ1 -- Specification quality.] Does ReversaFeynman preserve or improve the coverage, traceability, and correctness of operational specifications relative to the original Reversa pipeline?
  \item[RQ2 -- Repair effectiveness.] Does structural localization plus multi-candidate repair and independent validation improve issue-resolution success and reduce regressions relative to simpler repair baselines?
  \item[RQ3 -- Epistemic reliability.] Does explicit evidence governance reduce the rate at which inferred or unsupported claims are represented as directly observed?
  \item[RQ4 -- Calibration.] Are route and success-confidence estimates better calibrated after offline adaptive evaluation?
  \item[RQ5 -- Efficiency.] What are the latency, token, execution, and monetary cost trade-offs introduced by the additional governance and validation layers?
  \item[RQ6 -- Optimization safety.] Can shadow-first prompt/skill optimization improve measured task outcomes without increasing regression or epistemic-error rates?
\end{description}

\subsection{Systems under comparison}

A controlled evaluation should compare at least four variants:

\begin{enumerate}
  \item \textbf{Reversa-original}: the original reverse documentation workflow as the provenance baseline;
  \item \textbf{ReversaFeynman-core}: Reversa plus epistemic governance, without the v5 repair intelligence;
  \item \textbf{Minimal-repair baseline}: a simple localization/repair/validation loop inspired by the minimal-agent and Agentless literature \cite{xia2024agentless};
  \item \textbf{ReversaFeynman-v5}: the complete evidence-governed software engineering pipeline.
\end{enumerate}

All variants should use the same underlying model, temperature, repository snapshot, task statement, tool budget, timeout, and retry policy within each paired comparison.

\subsection{Task suites}

The evaluation should combine three task families:

\begin{itemize}
  \item \textbf{Real-world issue repair}: public benchmark instances from SWE-bench-style repositories \cite{jimenez2024swebench};
  \item \textbf{Generated repository tasks}: reproducible tasks generated with SWE-smith-style environment construction and test-breaking procedures \cite{yang2025swesmith};
  \item \textbf{Legacy specification tasks}: repositories selected for rule recovery, architecture reconstruction, specification generation, and migration planning, including cases where the original Reversa pipeline can be reproduced.
\end{itemize}

Repository-level splits must prevent the same repository from leaking across development/tuning and final evaluation when learned policies or prompt optimizers are involved.

\subsection{Primary metrics}

\paragraph{Specification metrics.}
We propose claim coverage, traceability coverage, human-adjudicated claim precision, unresolved-gap preservation, and contradiction rate. A separate \emph{false-observed rate} measures the fraction of claims labeled \texttt{OBSERVED} that lack sufficient direct evidence under independent review.

\paragraph{Repair metrics.}
Primary measures are issue resolution rate, pass@1, pass@k for candidate patches, regression rate, test-suite delta, mutation score, and the proportion of repairs for which localization metadata is traceable to the final patch.

\paragraph{Calibration metrics.}
When a system emits a probability of task success, calibration is evaluated using Brier score \cite{brier1950verification}, Expected Calibration Error (ECE), reliability diagrams, and coverage conditional on confidence bins. Calibration is computed only where a binary ground-truth outcome exists.

\paragraph{Efficiency metrics.}
We measure end-to-end latency, model tokens, tool calls, execution count, candidate patches evaluated, estimated monetary cost, and human interventions.

\paragraph{Governance metrics.}
We measure abstention precision, drift-trigger frequency, blocked-state resolution rate, evidence-promotion rejection rate, and the number of policy/skill proposals that remain shadow versus become eligible for activation.

\subsection{Experimental design and statistics}

For issue-repair tasks, systems should be evaluated in a paired design over the same task instances. When stochasticity is material, each system-task pair should be repeated with multiple independent seeds. Repository should be treated as a clustering unit rather than assuming all tasks are independent.

We recommend reporting effect sizes with 95\% confidence intervals using a repository-aware or hierarchical bootstrap. Paired binary outcomes can additionally be examined with McNemar's test; continuous paired metrics may use permutation tests or Wilcoxon signed-rank tests when parametric assumptions are not justified. When many hypotheses are tested, family-wise error or false-discovery control should be pre-specified.

For adaptive-policy evaluation, chronological or repository-separated holdouts must be maintained. Offline policy metrics should not be interpreted as causal effects unless propensities, overlap assumptions, and an appropriate off-policy estimator are available. A promotion-readiness decision is therefore an engineering gate, not a scientific claim of causal superiority.

\subsection{Reproducibility package}

Each reported experiment should preserve:

\begin{itemize}
  \item repository URL and immutable commit SHA;
  \item task identifier and test harness version;
  \item model/provider/version and decoding parameters;
  \item ReversaFeynman commit/tag;
  \item configuration of adapters and quality thresholds;
  \item random seeds;
  \item execution traces and normalized results;
  \item ReversaBench records and analysis scripts;
  \item environment/container identifiers when applicable.
\end{itemize}

This package is necessary to convert architectural claims into independently testable empirical claims.
\section{Experimental Status and Result-Gating}

The evaluation protocol in the previous section is now backed by an executable \emph{ReversaBench Experimental Harness}. The harness implements versioned task/run/report contracts, task--seed pairing, deterministic percentile bootstrap confidence intervals, regression and cost aggregation, and an explicit \emph{false-observed rate}. It also preserves ablation metadata for the Feynman layer, Hermes Evidence Governor, Code Intelligence, multi-candidate repair, mutation gate, and offline optimizer.

\subsection{Engineering smoke validation versus scientific evidence}

The repository contains a deterministic smoke manifest whose only purpose is to verify the benchmark machinery. Smoke records are tagged with \texttt{smoke=true}. When a caller requests a confirmatory report and only smoke records are available, the harness returns \texttt{confirmatory=false} with a blocking reason. Synthetic smoke outcomes therefore cannot be promoted into a confirmatory table by configuration alone.

This distinction is deliberate. Smoke execution can validate serialization, pairing, aggregation, confidence-interval generation, report rendering, and CI plumbing, but it does not test whether ReversaFeynman performs better than Reversa-original or another coding agent.

\subsection{Confirmatory result gate}

A confirmatory ReversaBench report requires non-smoke task runs with immutable repository and commit provenance. Each run retains task identifier, variant, seed, latency, cost, regression count, observed-claim count, false-observed-claim count, and ablation state. The default report continues to declare \texttt{causal\_claim=false}; causality is a property of the experimental design, not of the aggregation code.

The manuscript may include generated tables when a result artifact exists:

\IfFileExists{generated/reversabench-confirmatory.tex}{%
  \input{generated/reversabench-confirmatory}%
}{%
  \paragraph{Current result status.} No confirmatory ReversaBench result artifact is committed in this version. Consequently, the paper makes no empirical superiority claim. The repository contains the executable harness and pre-specified protocol needed to produce such a result under controlled runs.
}

\subsection{Artifact path}

The experimental implementation is linked to the paper through the following artifacts:
\begin{itemize}
  \item \texttt{specs/SPEC-REVERSABENCH-EXPERIMENTAL-HARNESS-V1.md};
  \item \texttt{lib/integrations/software-engineering/experimental-harness.js};
  \item \texttt{scripts/test-reversabench-experimental-harness.mjs};
  \item \texttt{benchmarks/reversabench/manifest.smoke.json};
  \item \texttt{scripts/run-reversabench-smoke.mjs};
  \item \texttt{scripts/render-reversabench-report.mjs}.
\end{itemize}

This gating mechanism ensures that the distinction between an implemented evaluation instrument and a completed empirical study remains machine-enforced as well as editorially stated.

\section{Related Work}

\subsection{Reverse documentation and legacy understanding}

Reversa is the direct foundation of this work. It frames legacy analysis as reverse documentation engineering and emphasizes traceability, explicit confidence marking, preserved gaps, and operational specifications for downstream agents \cite{macedo2026reversa}. ReversaFeynman retains this problem definition and extends the system after extraction with evidence governance and software engineering execution layers.

\subsection{Automated software engineering and repair}

Agentless demonstrated that a comparatively simple localization--repair--validation pipeline can be competitive with more complex autonomous agent scaffolds on software-engineering benchmarks \cite{xia2024agentless}. This result motivates ReversaFeynman's separation of structural localization, multiple patch generation, and independent validation rather than assuming that increasingly complex orchestration is always beneficial.

SWE-bench established a widely used evaluation setting for real-world GitHub issue resolution \cite{jimenez2024swebench}. SWE-smith later addressed the scarcity and cost of training/evaluation data by generating large numbers of software engineering tasks and execution environments \cite{yang2025swesmith}. ReversaBench is designed to interoperate with these styles of evaluation while adding specification and epistemic-governance metrics.

\subsection{Repository context and static program analysis}

Aider's repository-map approach uses Tree-sitter-derived structural information to summarize large repositories for LLM context selection \cite{aider2026repomap}. ReversaFeynman's Code Intelligence Graph generalizes this idea behind a provider-neutral contract. Semantic static analysis systems such as Semgrep provide another source of structural and policy-oriented findings \cite{semgrep2026}. These findings are treated as validation signals rather than automatic truth claims.

\subsection{Mutation testing and falsifiability}

Mutation testing evaluates the adequacy of tests by introducing controlled program changes and measuring whether the test suite detects them. ReversaFeynman uses mutation-testing results, for example from StrykerJS-style providers, as an executable probe for FEG-04 \cite{stryker2026}. A high mutation score indicates stronger sensitivity of the current oracle to the tested mutation operators; it does not by itself establish functional correctness.

\subsection{Agent observability and durable orchestration}

Modern agent systems increasingly require traces, experiments, and persistent state. Phoenix provides open-source observability and evaluation facilities for LLM applications \cite{phoenix2026}, while LangGraph provides stateful and durable orchestration primitives \cite{langgraph2026}. ReversaFeynman integrates these concerns through neutral tracing and checkpoint interfaces so that observability and workflow engines do not own epistemic policy.

\subsection{Programmatic prompt and agent optimization}

DSPy treats LM pipelines as programmable and optimizable modules rather than fixed prompt strings \cite{khattab2024dspy}. GEPA extends this direction with reflective prompt evolution over trajectories and demonstrates that language-level optimization can compete with or outperform some reinforcement-learning approaches while using fewer rollouts in the reported tasks \cite{agrawal2026gepa}. ReversaFeynman permits such systems as offline optimization providers, but optimization results remain shadow proposals until independent evaluation and governance gates are satisfied.

\subsection{Positioning}

The distinguishing goal of ReversaFeynman is not to replace specialized systems for execution, static analysis, observability, workflow durability, or optimization. Instead, it provides an evidence-governed integration architecture in which those systems remain optional and their scores do not silently acquire epistemic authority.
\section{Threats to Validity}

\paragraph{Construct validity.}
Terms such as \emph{understanding}, \emph{evidence}, and \emph{epistemic error} can be interpreted differently across studies. We mitigate this by operationalizing states and gates in versioned contracts, but the mapping between these engineering constructs and broader philosophical notions remains limited. Likewise, mutation score measures test sensitivity to selected mutation operators rather than complete behavioral correctness.

\paragraph{Internal validity.}
If future experiments compare systems using different models, tool budgets, temperatures, retries, or repository snapshots, observed differences may be attributable to those factors rather than to the ReversaFeynman architecture. Paired task execution with controlled configurations and repeated seeds is therefore necessary. Adaptive components also create temporal leakage risks; the explicit \texttt{decide()}--\texttt{observe()} separation and holdout requirements are intended to reduce this risk but do not eliminate all forms of contamination.

\paragraph{External validity.}
SWE-bench-style issue repair does not represent all maintenance work, and generated SWE-smith-style tasks may differ from naturally occurring defects. Legacy business systems can contain domain constraints that are absent from public open-source repositories. Evaluation should therefore span issue repair, specification recovery, migration planning, and multiple languages/domains before broad generalization.

\paragraph{Conclusion validity.}
Agent outcomes can be highly stochastic and clustered by repository. Treating all issue instances as independent can produce overly narrow confidence intervals. We therefore recommend repository-aware resampling, paired comparisons, multiple seeds, effect sizes, and pre-specified statistical tests. Offline policy comparisons should not be interpreted causally without suitable logging and identification assumptions.

\paragraph{Implementation validity.}
Several external capabilities are represented through adapters rather than embedded dependencies. The quality of a Tree-sitter extractor, execution provider, static analyzer, mutation framework, or observability backend can materially affect results. Experiments must record provider versions and configurations. A provider failure should remain distinguishable from a framework-level failure.

\paragraph{Evolution bias.}
ReversaFeynman was designed by iteratively studying weaknesses in earlier versions and patterns from related tools. This creates a risk of overfitting the architecture to known benchmark styles or favored engineering assumptions. The inclusion of minimal baselines, unseen repositories, and ablation studies is essential to determine which layers provide measurable value and which merely add complexity.

\paragraph{Current empirical status.}
The present manuscript describes an implemented architecture and evaluation protocol. It does not yet provide a controlled benchmark demonstrating superiority over the original Reversa or other software engineering agents. Any stronger claim would exceed the available evidence.
\section{Conclusion}

ReversaFeynman extends reverse documentation engineering into an evidence-governed software engineering architecture for agentic maintenance. The original Reversa contribution remains the foundation for extracting traceable operational specifications from legacy systems. ReversaFeynman adds explicit epistemic states, falsifiability-oriented gates, governed human teach-back, adaptive policy evaluation, Hermes-based memory and evidence collection, deterministic evidence governance, and a Software Engineering Intelligence layer for structural localization, isolated execution, multi-candidate repair, quality gates, observability, benchmarking, durable workflows, and shadow-first optimization.

The architectural thesis is that adaptive systems should be allowed to become better at choosing actions without being allowed to redefine what counts as direct evidence. In practical terms, memory, confidence, reward, benchmark performance, static-analysis scores, mutation scores, and optimizer rankings can influence routing and prioritization, but they remain distinct from traceable execution evidence.

The implementation is accompanied by versioned specifications, regression tests, strict contracts, and a proposed experimental protocol. The next research step is empirical rather than architectural: populate ReversaBench with reproducible legacy-specification and issue-repair tasks, run controlled ablations against the original Reversa and minimal repair baselines, and quantify whether the additional layers improve specification reliability, repair success, calibration, regression avoidance, and cost effectiveness.

Accordingly, the current contribution should be read as an implemented research framework and a falsifiable evaluation plan. A claim of superiority is intentionally deferred until the proposed experiments produce sufficient evidence.

\bibliographystyle{plain}
\bibliography{references}

\end{document}
